%%%%%%%%%%%%%%%%%%%%%%%%%%%%%%%%%%%%%%%%%%%%%%%%%%%%%%%%%%%%%%%%%%%%%%%%%%%%%%%
% APPENDIX AI: Spatial and Physical Context (Ψ_P)
% Geography, Location, Cities, and Spatial Economics
%%%%%%%%%%%%%%%%%%%%%%%%%%%%%%%%%%%%%%%%%%%%%%%%%%%%%%%%%%%%%%%%%%%%%%%%%%%%%%%
\section{Spatial and Physical Context: Geography and Location}
\label{app:spatial-context}

%==============================================================================
% PART 0: INTEGRATION OVERVIEW
%==============================================================================
\subsection{Framework Integration Map}

This appendix fills the final gap in the $(C, \Psi, C^*, K, Q)$ framework: 
the physical/spatial dimension $\Psi_P$. Where economic activity occurs 
profoundly shapes outcomes---agglomeration economies, transportation costs, 
natural resources, climate, and urban structure all constitute physical 
context. The New Economic Geography (Krugman) and Urban Economics (Glaeser) 
show how spatial $C^*$ emerges from the interaction of increasing returns, 
transportation costs, and factor mobility.

\begin{center}
\renewcommand{\arraystretch}{1.4}
\begin{tabular}{|l|l|l|}
\hline
\textbf{Framework Element} & \textbf{Spatial Contribution} & \textbf{Key Papers} \\
\hline
$\Psi_P$ (location) & Geography shapes opportunity & Krugman 1991 \\
Agglomeration & $C_{ij} > 0$ in space & Marshall 1890, Glaeser 2008 \\
Spatial $C^*$ & Core-periphery equilibria & Fujita et al. 1999 \\
Path dependence & Location lock-in & Davis \& Weinstein 2002 \\
Urban premium & Cities and productivity & Glaeser \& Gottlieb 2009 \\
Place-based policy & $\Psi_{design}$ for location & Kline \& Moretti 2014 \\
\hline
\end{tabular}
\end{center}

\paragraph{Why $\Psi_P$ Matters.}
Physical context shapes economic outcomes fundamentally:
\begin{enumerate}
    \item \textbf{Agglomeration}: Productivity rises with density
    \item \textbf{Market access}: Distance determines trade costs
    \item \textbf{Natural resources}: Geography as endowment
    \item \textbf{Climate}: Physical environment affects productivity
    \item \textbf{Infrastructure}: Transportation networks shape $C^*$
    \item \textbf{Sorting}: People and firms choose locations
\end{enumerate}

%==============================================================================
% PART I: FOUNDATIONS – CLASSICAL LOCATION THEORY
%==============================================================================
\subsection{Part I: Foundations – Classical Location Theory}

\subsubsection{von Thünen: Agricultural Location}

Johann Heinrich von Thünen (1826) pioneered spatial economics:

\begin{equation}
\boxed{
\text{Land rent} = \text{Yield} \times (\text{Price} - \text{Transport cost} \times \text{Distance})
}
\label{eq:thunen}
\end{equation}

\paragraph{The Model.}
\begin{itemize}
    \item Isolated city surrounded by agricultural land
    \item Different crops have different transport costs
    \item Concentric rings of land use emerge
\end{itemize}

\subsubsection{Marshall: Agglomeration Economies}

Alfred Marshall (1890) identified three sources of agglomeration:

\begin{equation}
\boxed{
\text{Agglomeration} = \text{Labor pooling} + \text{Input sharing} + \text{Knowledge spillovers}
}
\label{eq:marshall-agglomeration}
\end{equation}

\paragraph{Marshallian Externalities.}
\begin{enumerate}
    \item \textbf{Labor market pooling}: Thick markets reduce matching frictions
    \item \textbf{Input sharing}: Specialized suppliers viable at scale
    \item \textbf{Knowledge spillovers}: Ideas flow locally
\end{enumerate}

\subsubsection{Framework Translation}

Agglomeration is spatial complementarity:
\begin{equation}
C_{ij}^{spatial} > 0 \quad \text{(co-location is complementary)}
\end{equation}

\subsubsection{Weber and Lösch: Industrial Location}

Weber (1909) and Lösch (1940) developed industrial location theory:
\begin{itemize}
    \item Firms locate to minimize transport costs
    \item Market areas form hexagonal patterns
    \item Central place hierarchy emerges
\end{itemize}

%==============================================================================
% PART II: NEW ECONOMIC GEOGRAPHY
%==============================================================================
\subsection{Part II: New Economic Geography – Krugman's Revolution}

\subsubsection{Krugman: Geography and Trade}

Paul Krugman (1991) launched the New Economic Geography:

\begin{equation}
\boxed{
\text{Spatial } C^* = f(\text{Increasing returns}, \text{Transport costs}, \text{Factor mobility})
}
\label{eq:neg}
\end{equation}

\paragraph{The Core-Periphery Model.}
\begin{itemize}
    \item Two regions, manufacturing and agriculture
    \item Manufacturing has increasing returns
    \item Workers mobile, farmers immobile
    \item Transport costs for manufactures
\end{itemize}

\subsubsection{Multiple Equilibria}

The model generates multiple spatial equilibria:

\begin{center}
\renewcommand{\arraystretch}{1.3}
\begin{tabular}{|l|l|}
\hline
\textbf{Transport Costs} & \textbf{Equilibrium} \\
\hline
Very high & Dispersed (autarky) \\
Intermediate & Core-periphery (agglomeration) \\
Very low & Dispersed (factor price equalization) \\
\hline
\end{tabular}
\end{center}

\subsubsection{Framework Translation}

Spatial $C^*$ depends on $\Psi_P$:
\begin{equation}
C^*_{spatial} \in \{C^*_{dispersed}, C^*_{core-periphery}\}
\end{equation}

Which equilibrium depends on transport costs and increasing returns.

\subsubsection{Centripetal vs. Centrifugal Forces}

\begin{center}
\renewcommand{\arraystretch}{1.3}
\begin{tabular}{|l|l|}
\hline
\textbf{Centripetal (Agglomeration)} & \textbf{Centrifugal (Dispersion)} \\
\hline
Market size effects & Immobile factors \\
Thick labor markets & Land/housing costs \\
Knowledge spillovers & Congestion \\
Input-output linkages & Competition for workers \\
\hline
\end{tabular}
\end{center}

\subsubsection{Fujita, Krugman \& Venables: The Spatial Economy}

Fujita, Krugman \& Venables (1999) provided comprehensive treatment:

\begin{equation}
\text{Urban system} = \text{Hierarchy of cities with different functions}
\end{equation}

\paragraph{Key Results.}
\begin{itemize}
    \item City size distribution (Zipf's law)
    \item Urban hierarchy emerges endogenously
    \item Specialization across cities
\end{itemize}

%==============================================================================
% PART III: URBAN ECONOMICS
%==============================================================================
\subsection{Part III: Urban Economics – Cities as Engines}

\subsubsection{Glaeser: Triumph of the City}

Edward Glaeser emphasized cities as drivers of growth:

\begin{equation}
\boxed{
\text{Cities} = \text{Absence of space between people and firms}
}
\label{eq:glaeser-city}
\end{equation}

\paragraph{Why Cities Exist.}
\begin{itemize}
    \item Reduce transport costs for goods, people, ideas
    \item Enable specialization and matching
    \item Facilitate knowledge spillovers
\end{itemize}

\subsubsection{Urban Wage Premium}

Glaeser \& Maré (2001) documented:

\begin{equation}
\boxed{
\text{Urban wage premium} \approx 25\text{--}35\%
}
\label{eq:urban-premium}
\end{equation}

\paragraph{Sources.}
\begin{itemize}
    \item Productivity (agglomeration economies)
    \item Sorting (skilled workers move to cities)
    \item Learning (faster human capital accumulation)
\end{itemize}

\subsubsection{Framework Translation}

Cities are high-$\Psi_P$ environments:
\begin{equation}
\Psi_P^{city} > \Psi_P^{rural} \quad \text{(more opportunities, interactions)}
\end{equation}

\subsubsection{Agglomeration Elasticities}

Rosenthal \& Strange (2004) surveyed estimates:

\begin{equation}
\frac{\partial \ln(\text{productivity})}{\partial \ln(\text{density})} \approx 0.03\text{--}0.08
\end{equation}

Doubling density raises productivity 3--8\%.

\subsubsection{Combes \& Gobillon: Agglomeration Survey}

Combes \& Gobillon (2015) provided comprehensive review:
\begin{itemize}
    \item Agglomeration effects are real and substantial
    \item Sorting explains part but not all of urban premium
    \item Mechanisms: matching, sharing, learning
\end{itemize}

%==============================================================================
% PART IV: GEOGRAPHY AND DEVELOPMENT
%==============================================================================
\subsection{Part IV: Geography and Development}

\subsubsection{Sachs: Geography Matters}

Jeffrey Sachs argued geography directly affects development:

\begin{equation}
\boxed{
\text{Tropics, landlocked, malaria} \Rightarrow \text{Lower income}
}
\label{eq:sachs-geography}
\end{equation}

\paragraph{Direct Effects.}
\begin{itemize}
    \item Disease burden (malaria, tropical diseases)
    \item Agricultural productivity (soil, climate)
    \item Market access (landlocked, distance)
\end{itemize}

\subsubsection{Gallup, Sachs \& Mellinger}

Gallup, Sachs \& Mellinger (1999) documented:
\begin{itemize}
    \item Coastal regions much richer than interior
    \item Temperate zones richer than tropics
    \item Market access strongly predicts income
\end{itemize}

\subsubsection{Diamond: Guns, Germs, and Steel}

Jared Diamond (1997) argued ultimate causes are geographic:

\begin{equation}
\text{Biogeography} \to \text{Agriculture} \to \text{States} \to \text{Technology}
\end{equation}

\paragraph{Key Argument.}
\begin{itemize}
    \item Eurasia had more domesticable plants/animals
    \item East-west axis allowed diffusion
    \item Geography determined which societies developed first
\end{itemize}

\subsubsection{Acemoglu et al.: Institutions vs. Geography}

Acemoglu, Johnson \& Robinson (2002) challenged geography:

\begin{equation}
\text{Institutions} \gg \text{Geography} \quad \text{(for development)}
\end{equation}

\paragraph{The Debate.}
\begin{itemize}
    \item Geography may affect institutions (colonial origins)
    \item Direct geography effects small once institutions controlled
    \item Reversal of fortune: rich tropical colonies now poor
\end{itemize}

\subsubsection{Framework Translation}

Geography affects development through multiple channels:
\begin{equation}
\Psi_P \to \Psi_I \to C^* \to Q \quad \text{and} \quad \Psi_P \to Q \text{ directly}
\end{equation}

%==============================================================================
% PART V: SPATIAL EQUILIBRIUM AND SORTING
%==============================================================================
\subsection{Part V: Spatial Equilibrium and Sorting}

\subsubsection{Rosen-Roback Framework}

Rosen (1979) and Roback (1982) developed spatial equilibrium:

\begin{equation}
\boxed{
\text{Utility equalized across locations: } V(w_j, r_j, A_j) = \bar{V}
}
\label{eq:roback}
\end{equation}

where $w$ = wage, $r$ = rent, $A$ = amenities.

\paragraph{Implications.}
\begin{itemize}
    \item High amenity places have high rents and/or low wages
    \item Wages compensate for disamenities
    \item Housing prices capitalize amenities
\end{itemize}

\subsubsection{Framework Translation}

Spatial $C^*$ equalizes utility:
\begin{equation}
C^*_{spatial}: \quad Q_j = Q_k \quad \forall j, k \text{ (for mobile households)}
\end{equation}

\subsubsection{Moretti: Geography of Jobs}

Enrico Moretti (2012) showed:

\begin{equation}
\text{Great Divergence} = \text{Innovation hubs vs. declining cities}
\end{equation}

\paragraph{Key Findings.}
\begin{itemize}
    \item College graduates increasingly concentrate
    \item Innovation jobs have large multipliers (5x)
    \item Place-based divergence in wages and opportunities
\end{itemize}

\subsubsection{Autor et al.: Place-Based Inequality}

Autor (2019) documented:
\begin{itemize}
    \item Urban wage premium rising over time
    \item Non-college workers not benefiting from cities
    \item Spatial sorting by skill intensifying
\end{itemize}

%==============================================================================
% PART VI: PATH DEPENDENCE IN SPACE
%==============================================================================
\subsection{Part VI: Path Dependence – Location Lock-In}

\subsubsection{Davis \& Weinstein: Bones of Bombed Cities}

Davis \& Weinstein (2002) studied Japanese cities after WWII:

\begin{equation}
\boxed{
\text{Cities returned to pre-war size within 15--20 years}
}
\label{eq:davis-weinstein}
\end{equation}

\paragraph{Interpretation.}
\begin{itemize}
    \item Fundamentals (geography) determine city location
    \item Even massive destruction doesn't relocate cities
    \item Multiple equilibria exist but fundamentals select
\end{itemize}

\subsubsection{Bleakley \& Lin: Portage Sites}

Bleakley \& Lin (2012) studied portage sites:
\begin{itemize}
    \item Historical portage sites became cities
    \item Portage no longer relevant (canals, railroads)
    \item Cities persist due to sunk costs, agglomeration
\end{itemize}

\paragraph{Framework Translation.}
Spatial path dependence:
\begin{equation}
C^*_{spatial,t} = f(\Psi_P^{historical}, C^*_{spatial,t-1})
\end{equation}

Historical geography shapes current economic geography.

\subsubsection{Redding et al.: Quantitative Spatial Economics}

Redding \& Rossi-Hansberg (2017) surveyed quantitative models:
\begin{itemize}
    \item Tractable spatial general equilibrium models
    \item Counterfactual analysis of spatial policies
    \item Trade, migration, and spatial structure
\end{itemize}

%==============================================================================
% PART VII: PLACE-BASED POLICIES
%==============================================================================
\subsection{Part VII: Place-Based Policies – Designing $\Psi_P$}

\subsubsection{Kline \& Moretti: Place-Based Policies}

Kline \& Moretti (2014) evaluated place-based policies:

\begin{equation}
\boxed{
\text{Place-based policy} = \text{Subsidize locations, not people}
}
\label{eq:place-based}
\end{equation}

\paragraph{Arguments For.}
\begin{itemize}
    \item Multiple equilibria: can shift to better $C^*$
    \item Agglomeration externalities: private location choice inefficient
    \item Equity: help people who can't move
\end{itemize}

\paragraph{Arguments Against.}
\begin{itemize}
    \item May be inefficient (subsidize wrong places)
    \item Capital mobile, benefits may not stay local
    \item Better to help people directly
\end{itemize}

\subsubsection{Busso et al.: Enterprise Zones}

Busso, Gregory \& Kline (2013) evaluated Empowerment Zones:
\begin{itemize}
    \item Federal program for distressed areas
    \item Significant employment increases
    \item Wages rose, poverty fell
\end{itemize}

\subsubsection{Opportunity Zones}

Opportunity Zones (2017 tax law):
\begin{itemize}
    \item Capital gains deferral for investment in zones
    \item Early evidence mixed
    \item Design matters for effectiveness
\end{itemize}

\subsubsection{Framework Translation}

Place-based policy is $\Psi_P$ design:
\begin{equation}
\Psi_P^{policy} \to C^*_{spatial}^{new} \quad \text{(shift spatial equilibrium)}
\end{equation}

%==============================================================================
% PART VIII: SYNTHESIS
%==============================================================================
\subsection{Part VIII: Synthesis – Spatial Context Equations}

\subsubsection{Complete Framework Translation}

\begin{center}
\renewcommand{\arraystretch}{1.5}
\begin{tabular}{|l|c|l|}
\hline
\textbf{Concept} & \textbf{Equation} & \textbf{$\Psi_P$ Element} \\
\hline
Transport costs & $\tau_{ij}$ & Distance friction \\
\hline
Agglomeration & $A(L) = L^\gamma$ & Density externality \\
\hline
Core-periphery & Multiple $C^*_{spatial}$ & Spatial equilibrium \\
\hline
Urban premium & $w_{city} / w_{rural} \approx 1.3$ & Location productivity \\
\hline
Spatial sorting & $V(w,r,A) = \bar{V}$ & Equilibrium condition \\
\hline
Path dependence & $C^*_t = f(C^*_{t-1}, \Psi_P)$ & History matters \\
\hline
Place-based policy & $\Psi_P^{design}$ & Spatial intervention \\
\hline
\end{tabular}
\end{center}

\subsubsection{$\Psi_P$ in the Full Framework}

\begin{equation}
\boxed{
C^* = C^*(\Psi_P, \Psi_C, \Psi_I, \Psi_T, \ldots)
}
\end{equation}

Physical context interacts with other dimensions:
\begin{itemize}
    \item $\Psi_P \times \Psi_I$: Geography shapes institutions
    \item $\Psi_P \times \Psi_K$: Knowledge spillovers are local
    \item $\Psi_P \times \Psi_S$: Social networks are spatial
    \item $\Psi_P \times \Psi_E$: Housing prices vary spatially
\end{itemize}

\subsubsection{Welfare Implications}

Spatial misallocation reduces welfare:
\begin{equation}
Q^{actual} < Q^{optimal} \quad \text{due to spatial frictions}
\end{equation}

Hsieh \& Moretti (2019) estimate US GDP 36\% lower due to housing constraints 
preventing workers from moving to productive cities.

%==============================================================================
% PART IX: COMPLETE PAPER DOCUMENTATION
%==============================================================================
\subsection{Part IX: Complete Paper Documentation}

%--- CLASSICAL ---
\subsubsection{Classical Location Theory}

\paragraph{Paper 1: von Thünen (1826)}

\textit{Der isolierte Staat in Beziehung auf Landwirtschaft und Nationalökonomie.}

\textbf{Summary.}
First spatial economic model. Agricultural land use varies with distance 
from city. Concentric rings of different crops.

\textbf{Framework Translation.}
$\Psi_P$ (distance to market) determines land use $C^*$.

\paragraph{Paper 2: Marshall (1890)}

\textit{Principles of Economics.}
Macmillan.

\textbf{Summary.}
Identifies three sources of agglomeration: labor pooling, input sharing, 
knowledge spillovers. ``Mysteries of the trade become no mystery.''

\textbf{Framework Translation.}
\begin{equation}
C_{ij}^{spatial} > 0 \quad \text{(Marshallian externalities)}
\end{equation}

\textbf{Key Finding.}
Foundation of agglomeration economics.

%--- NEW ECONOMIC GEOGRAPHY ---
\subsubsection{New Economic Geography}

\paragraph{Paper 3: Krugman (1991)}

``Increasing Returns and Economic Geography.''
\textit{Journal of Political Economy}, 99(3), 483--499.

\textbf{Summary.}
Launches New Economic Geography. Core-periphery model with increasing 
returns, transport costs, and labor mobility. Multiple equilibria.

\textbf{Framework Translation.}
\begin{equation}
C^*_{spatial} \in \{\text{dispersed}, \text{agglomerated}\}
\end{equation}

\textbf{Key Finding.}
Nobel Prize 2008. 10,000+ citations.

\paragraph{Paper 4: Krugman (1991b)}

\textit{Geography and Trade.}
MIT Press.

\textbf{Summary.}
Book treatment of economic geography. Trade patterns, location of 
production, regional economics.

\textbf{Framework Translation.}
Complete spatial economics framework.

\paragraph{Paper 5: Fujita, Krugman \& Venables (1999)}

\textit{The Spatial Economy: Cities, Regions, and International Trade.}
MIT Press.

\textbf{Summary.}
Definitive treatment of New Economic Geography. Urban systems, regional 
structures, international trade patterns.

\textbf{Framework Translation.}
Comprehensive theory of spatial $C^*$.

\textbf{Key Finding.}
Definitive NEG textbook. 8,000+ citations.

%--- URBAN ECONOMICS ---
\subsubsection{Urban Economics}

\paragraph{Paper 6: Glaeser (2008)}

\textit{Cities, Agglomeration, and Spatial Equilibrium.}
Oxford University Press.

\textbf{Summary.}
Modern urban economics. Spatial equilibrium, agglomeration, sorting, 
amenities, housing.

\textbf{Framework Translation.}
$\Psi_P^{urban}$ shapes economic outcomes.

\paragraph{Paper 7: Glaeser (2011)}

\textit{Triumph of the City: How Our Greatest Invention Makes Us Richer, 
Smarter, Greener, Healthier, and Happier.}
Penguin.

\textbf{Summary.}
Popular treatment of urban economics. Cities as engines of innovation, 
environmentally friendly, key to prosperity.

\textbf{Framework Translation.}
Cities maximize $Q$ through agglomeration.

\paragraph{Paper 8: Glaeser \& Gottlieb (2009)}

``The Wealth of Cities: Agglomeration Economies and Spatial Equilibrium.''
\textit{Journal of Economic Literature}, 47(4), 983--1028.

\textbf{Summary.}
Survey of urban economics. Agglomeration, spatial equilibrium, 
urban growth, policy implications.

\textbf{Framework Translation.}
State of the art on cities and $\Psi_P$.

\paragraph{Paper 9: Rosenthal \& Strange (2004)}

``Evidence on the Nature and Sources of Agglomeration Economies.''
\textit{Handbook of Regional and Urban Economics}, Vol. 4.

\textbf{Summary.}
Survey of agglomeration evidence. Estimates elasticities, tests 
mechanisms, reviews methods.

\textbf{Framework Translation.}
$\partial \ln(Y) / \partial \ln(\text{density}) \approx 0.03$--$0.08$.

\paragraph{Paper 10: Combes \& Gobillon (2015)}

``The Empirics of Agglomeration Economies.''
\textit{Handbook of Regional and Urban Economics}, Vol. 5.

\textbf{Summary.}
Updated survey of agglomeration. Sorting vs. agglomeration, mechanisms, 
dynamic effects.

\textbf{Framework Translation.}
Comprehensive agglomeration evidence.

%--- GEOGRAPHY AND DEVELOPMENT ---
\subsubsection{Geography and Development}

\paragraph{Paper 11: Gallup, Sachs \& Mellinger (1999)}

``Geography and Economic Development.''
\textit{International Regional Science Review}, 22(2), 179--232.

\textbf{Summary.}
Documents geographic correlates of development. Coastal access, 
tropics, landlocked status all predict income.

\textbf{Framework Translation.}
$\Psi_P \to Q$ (direct geographic effects).

\paragraph{Paper 12: Diamond (1997)}

\textit{Guns, Germs, and Steel: The Fates of Human Societies.}
W.W. Norton.

\textbf{Summary.}
Geographic determinism thesis. Biogeography determined which societies 
developed agriculture, states, technology first.

\textbf{Framework Translation.}
$\Psi_P^{biogeography} \to \Psi_I \to Q$.

\textbf{Key Finding.}
Pulitzer Prize. Influential but controversial.

\paragraph{Paper 13: Acemoglu, Johnson \& Robinson (2002)}

``Reversal of Fortune: Geography and Institutions in the Making of the 
Modern World Income Distribution.''
\textit{Quarterly Journal of Economics}, 117(4), 1231--1294.

\textbf{Summary.}
Challenges geography thesis. Reversal of fortune: rich pre-colonial 
areas now poor. Institutions matter more than geography.

\textbf{Framework Translation.}
$\Psi_I \gg \Psi_P$ for development.

%--- SPATIAL EQUILIBRIUM ---
\subsubsection{Spatial Equilibrium}

\paragraph{Paper 14: Roback (1982)}

``Wages, Rents, and the Quality of Life.''
\textit{Journal of Political Economy}, 90(6), 1257--1278.

\textbf{Summary.}
Spatial equilibrium framework. Utility equalized across locations 
through wages and rents. Amenities capitalized in prices.

\textbf{Framework Translation.}
\begin{equation}
C^*_{spatial}: V(w_j, r_j, A_j) = \bar{V}
\end{equation}

\textbf{Key Finding.}
Foundation of spatial equilibrium models.

\paragraph{Paper 15: Moretti (2012)}

\textit{The New Geography of Jobs.}
Houghton Mifflin Harcourt.

\textbf{Summary.}
Great divergence in American cities. Innovation hubs thriving, 
manufacturing cities declining. Each innovation job creates 5 local jobs.

\textbf{Framework Translation.}
$\Psi_P$ varies dramatically across cities.

%--- PATH DEPENDENCE ---
\subsubsection{Path Dependence}

\paragraph{Paper 16: Davis \& Weinstein (2002)}

``Bones, Bombs, and Break Points: The Geography of Economic Activity.''
\textit{American Economic Review}, 92(5), 1269--1289.

\textbf{Summary.}
Studies Japanese cities after WWII bombing. Cities returned to 
pre-war size. Fundamentals determine location.

\textbf{Framework Translation.}
Geographic fundamentals select among $C^*$.

\textbf{Key Finding.}
Influential on path dependence debate.

\paragraph{Paper 17: Bleakley \& Lin (2012)}

``Portage and Path Dependence.''
\textit{Quarterly Journal of Economics}, 127(2), 587--644.

\textbf{Summary.}
Historical portage sites became cities. Portage obsolete but cities 
persist. Path dependence in spatial structure.

\textbf{Framework Translation.}
$C^*_{spatial,t}$ depends on historical $\Psi_P$.

%--- QUANTITATIVE SPATIAL ---
\subsubsection{Quantitative Spatial Economics}

\paragraph{Paper 18: Redding \& Rossi-Hansberg (2017)}

``Quantitative Spatial Economics.''
\textit{Annual Review of Economics}, 9, 21--58.

\textbf{Summary.}
Survey of quantitative spatial models. Trade, migration, spatial 
equilibrium. Methods for counterfactual analysis.

\textbf{Framework Translation.}
Modern toolkit for analyzing $\Psi_P$.

\paragraph{Paper 19: Hsieh \& Moretti (2019)}

``Housing Constraints and Spatial Misallocation.''
\textit{American Economic Journal: Macroeconomics}, 11(2), 1--39.

\textbf{Summary.}
Housing constraints prevent workers from moving to productive cities.
US GDP 36\% lower due to spatial misallocation.

\textbf{Framework Translation.}
\begin{equation}
Q^{actual} < Q^{optimal} \quad \text{(spatial misallocation)}
\end{equation}

\textbf{Key Finding.}
Major policy implications.

%--- PLACE-BASED ---
\subsubsection{Place-Based Policy}

\paragraph{Paper 20: Kline \& Moretti (2014)}

``Place Based Policies with Unemployment.''
\textit{American Economic Review}, 104(5), 13--18.

\textbf{Summary.}
Theory of place-based policies. Can be efficient with multiple 
equilibria, agglomeration, or unemployment.

\textbf{Framework Translation.}
$\Psi_P^{design}$ can shift spatial $C^*$.

%%%%%%%%%%%%%%%%%%%%%%%%%%%%%%%%%%%%%%%%%%%%%%%%%%%%%%%%%%%%%%%%%%%%%%%%%%%%%%%
% COHERENT REGIONAL CLUSTERS
%%%%%%%%%%%%%%%%%%%%%%%%%%%%%%%%%%%%%%%%%%%%%%%%%%%%%%%%%%%%%%%%%%%%%%%%%%%%%%%
\subsection{Coherent Regional Clusters: Examples of Spatial $C^*$}
\label{sec:coherent-clusters}

The theory of agglomeration economies and spatial complementarity finds its
clearest expression in \textbf{coherent regional clusters}---geographic areas
where the elements of the regional economy form self-reinforcing systems.

\subsubsection{The Coherent Region}

A region with high coherence ($K_{region} \approx 1$) exhibits:
\begin{itemize}
    \item \textbf{Specialized ecosystem:} Firms, workers, institutions aligned
          around a dominant industry or activity
    \item \textbf{Self-reinforcing advantages:} Success attracts more talent,
          capital, and activity, which generates more success
    \item \textbf{High barriers to replication:} The complementary system cannot
          be copied piecemeal
\end{itemize}

\subsubsection{Case Study: Silicon Valley}

Silicon Valley exemplifies regional complementarity:

\begin{center}
\renewcommand{\arraystretch}{1.3}
\begin{tabular}{|l|l|}
\hline
\textbf{Element} & \textbf{Complementary Effect} \\
\hline
Stanford/Berkeley & Talent pipeline, research spillovers \\
Venture capital ecosystem & Risk capital for startups \\
Dense labor market & Easy hiring, job mobility \\
Legal/accounting expertise & Specialized support services \\
Entrepreneurial culture & Norm of risk-taking, failure tolerance \\
Network effects & Connections, information sharing \\
\hline
\end{tabular}
\end{center}

\paragraph{Framework Translation.}
All elements are strategic complements:
\begin{equation}
\gamma_{talent,VC} > 0, \quad \gamma_{VC,culture} > 0, \quad
\gamma_{culture,network} > 0, \quad \gamma_{network,talent} > 0
\end{equation}

This creates a \textbf{stable attractor} that resists replication. Austin, Boston,
or Tel Aviv can build parts of the system, but the complete coherent configuration
is difficult to recreate.

\subsubsection{Case Study: Wall Street}

New York's financial district exhibits similar complementarity:
\begin{itemize}
    \item \textbf{Density of financial institutions:} Banks, hedge funds, asset managers
    \item \textbf{Specialized labor market:} Deep pool of financial talent
    \item \textbf{Information flows:} Face-to-face contact for tacit knowledge
    \item \textbf{Legal/regulatory infrastructure:} Courts, regulators, law firms
    \item \textbf{Supporting services:} Restaurants, transportation, housing
\end{itemize}

Despite technological change (electronic trading, remote work), the complementarity
persists---the whole exceeds the sum of parts.

\subsubsection{Case Study: Hollywood}

Los Angeles entertainment industry shows:
\begin{itemize}
    \item \textbf{Talent agglomeration:} Actors, directors, writers, technicians
    \item \textbf{Studio infrastructure:} Production facilities, post-production
    \item \textbf{Deal-making culture:} Agents, managers, lawyers
    \item \textbf{Screening/testing:} Audience proximity, market research
    \item \textbf{Network effects:} ``Who you know'' matters for project assembly
\end{itemize}

\subsubsection{Regional Incoherence}

When regional elements are \textit{misaligned}, the result is incoherence
($K_{region} < 1$), which manifests as:
\begin{itemize}
    \item \textbf{Declining regions:} Loss of key industry without replacement
    \item \textbf{Brain drain:} Talent leaves for more coherent ecosystems
    \item \textbf{Economic stagnation:} Investment flows elsewhere
    \item \textbf{Institutional decay:} Supporting infrastructure erodes
\end{itemize}

\paragraph{Examples.}
Detroit after auto industry decline, Rust Belt manufacturing cities,
regions dependent on a single extractive industry.

\paragraph{Policy Implication.}
Place-based policies must address the \textit{system} of complementarities,
not individual elements. Building a research park without the venture capital,
talent pipeline, and entrepreneurial culture will not replicate Silicon Valley.

%%%%%%%%%%%%%%%%%%%%%%%%%%%%%%%%%%%%%%%%%%%%%%%%%%%%%%%%%%%%%%%%%%%%%%%%%%%%%%%
\subsection{Citation and Verification Note}

\textbf{Nobel Prizes represented:}
\begin{itemize}
    \item Krugman (2008) ``for his analysis of trade patterns and location of economic activity''
\end{itemize}

Combined Google Scholar citations for the 20 papers: $>$80,000.

This appendix provides the theoretical and empirical foundations for the 
spatial/physical dimension $\Psi_P$ of the framework, showing how geography 
and location fundamentally shape economic outcomes.

%%%%%%%%%%%%%%%%%%%%%%%%%%%%%%%%%%%%%%%%%%%%%%%%%%%%%%%%%%%%%%%%%%%%%%%%%%%%%%%
